\section{Conclusion}
\label{sec:conclusion}

The increasing complexity of thermal-aware dynamic scheduling algorithms dictates that the underlying thermal simulation engines must be transiently accurate and physically sound. In this work, we addressed a critical artifact in the gem5 simulator where intermediate thermal nodes defaulted to absolute zero, causing severe unphysical cooling during early simulation epochs. By employing a comprehensive five-way validation framework encompassing closed-form analytical solutions, independent numerical solvers, and industry-standard SPICE simulations, we not only diagnosed the artifact but proved the efficacy of our patch. Our results demonstrated a convergence to within \SI{0.08}{\degreeCelsius} of the physical ground truth. We argue that such ``physical admissibility proofs'' should become standard practice in the computer architecture community to restore confidence in native thermal simulation infrastructures and reduce the reliance on ad-hoc external toolchains.
